% ============================================================================
% config.tex -- configuração desta lista (edite isto, não o main.tex)
%
% Tudo que muda de lista pra lista -- ou que é fixo mas fica fora do
% repositório -- vai aqui. O main.tex (ou o .tex renomeado da lista) só
% deve ter conteúdo: capa + questões.
% ============================================================================

% --- Ativos institucionais (logo + ícone, sempre os mesmos) -----------------
% Ficam num arquivo externo, fora do controle de versão -- o mesmo
% caminho serve pra toda lista, então normalmente você NÃO precisa
% mexer nessa linha depois de configurá-la uma vez.
%
% Ajuste o caminho abaixo para o arquivo .tex externo que contém a sua
% chamada de \IFCEativosInstitucionais com os caminhos reais das
% imagens, por exemplo:
%
%   % conteúdo de /home/santiago/algum-lugar/institucional.tex:
%   \IFCEativosInstitucionais
%     {\includegraphics[width=7.25cm]{/home/santiago/algum-lugar/logo-ifce-baturite.pdf}}
%     {\includegraphics[width=0.37cm]{/home/santiago/algum-lugar/icone-ifce-mini.pdf}}
%
% Se o arquivo não existir (como no próprio template, sem nenhuma lista
% real configurada ainda), caem placeholders genéricos no lugar --
% a compilação nunca quebra por causa disso.
\InputIfFileExists{/CAMINHO/EXTERNO/institucional.tex}{}{%
  \IFCEativosInstitucionais
    {\IFCEplaceholder[4.0cm]{logo da instituição}}
    {\textcolor{ifcegreen}{\rule{0.3cm}{0.3cm}}}
}

% --- Metadados desta lista (edite a cada lista nova) ------------------------
% Os valores abaixo são os do exemplo de demonstração deste próprio
% template (main.tex); numa lista real, troque todos os 8 campos.
\IFCEmetadados
  {IFCE Baturité}                          % instituição
  {Técnico Integrado em Comércio}           % curso/turma
  {Física II/Complementar}                  % disciplina
  {Lista de Exercícios}                     % tipo de conteúdo
  {Enem e UECE}                             % descrição do conteúdo
  {Vetores}                                 % assunto/tópico
  {Paulo Santiago}                          % professor
  {\IFCEplaceholder[2.0cm]{logo do curso}}  % logomarca do curso (troque por
                                             % \includegraphics{...} quando tiver a arte)
